% Historical source (removed by author request): references/_OLD_Übungsaufgaben/Übungsblatt 3 - Vektorechnung/Blatt3.tex:177
\begin{ExerciseBlock}
\ExerciseHeading{09}{Winkel zwischen Vektoren}{vector_algebra_vector_angles}
\begin{enumerate}[label=\alph*)]

\item Berechnen Sie den Winkel $\alpha$ zwischen den Vektoren
\begin{center}
	$\vec{a}=\left(\begin{array}{c}  1 \\  1 \\  1  \end{array}\right) \ $ und 
	$\ \vec{b}=\left(\begin{array}{c}  1 \\  1 \\  -1  \end{array}\right)$
\end{center}

\item Bestimmen Sie ein $r \in \R$, sodass die Vektoren $\vec{a}$ und $\vec{b}$ senkrecht aufeinander stehen.
\begin{center}
$\vec{a}=\left(\begin{array}{c}  1 \\  1 \\  1  \end{array}\right)$ und $\vec{b}=\left(\begin{array}{c}  1 \\  1 \\  r  \end{array}\right)$
\end{center}

\end{enumerate}

\end{ExerciseBlock}
\UseExerciseSolution{vector_algebra_vector_angles}
